\chapter{AGORA2 Representative Strains and Oral-Fluid Medium}
\label{app:agora}

Table~\ref{tab:agora_strains} lists the AGORA2 v2.01 genome-scale model chosen as
the representative for each guild in the L2/L3 metabolic sign prior of
Chapter~\ref{ch:dieckow}. The \emph{Other} guild has no representative, so guild
pairs involving it receive no FBA-based prior.

\begin{table}[htbp]
  \centering\small
  \begin{tabular}{@{}ll@{}}
    \toprule
    Guild & AGORA2 representative strain\\
    \midrule
    Actinobacteria      & \textit{Actinomyces naeslundii} Howell 279\\
    Bacilli             & \textit{Streptococcus gordonii} CH1\\
    Negativicutes       & \textit{Veillonella parvula} Te3 DSM~2008\\
    Bacteroidia         & \textit{Porphyromonas melaninogenica} ATCC~25845\\
    Fusobacteriia       & \textit{Fusobacterium nucleatum} ATCC~25586\\
    Clostridia          & \textit{Parvimonas micra} ATCC~33270\\
    Coriobacteriia      & \textit{Atopobium parvulum} DSM~20469\\
    Gammaproteobacteria & \textit{Haemophilus parainfluenzae} T3T1\\
    Betaproteobacteria  & \textit{Eikenella corrodens} ATCC~23834\\
    Other               & no AGORA representative (L3 absent)\\
    \bottomrule
  \end{tabular}
  \caption[AGORA2 representative strains]{AGORA2 representative strains per guild
  used for the pFBA (L2) and MICOM community-FBA (L3) layers.}
  \label{tab:agora_strains}
\end{table}

\paragraph{Oral-fluid medium.}
The shared medium was assembled in two steps. First, quantified metabolite
concentrations in unstimulated whole saliva were drawn from the literature and
converted to exchange bounds (mmol\,gDW$^{-1}$\,h$^{-1}$) assuming a biofilm density
of ${\sim}2$\,gDW\,L$^{-1}$ and a reference growth rate of $0.1$\,h$^{-1}$:
sugars (glucose, fructose, sucrose, lactate), all twenty amino acids plus cysteine
and ornithine, B-vitamins, trace metals, purines, and inorganic salts.
Concentrations were set at the blood-plasma scale (${\sim}100\times$ realistic
saliva) so that every model remains feasible; a realistic-concentration variant
produced non-specific competition (growth suppressed $>50\%$ in $81/90$ pairs) and
was not used. Second, AGORA2-specific cofactors absent from the saliva list were
added at minimal bounds where required for feasibility (protoheme/siroheme,
menaquinones, \textit{meso}-diaminopimelate, long-chain fatty acids,
4-hydroxybenzoate, adenosylcobalamin, glutathione, polyamines, cytidine). The four
strictly anaerobic guilds had their O$_2$ exchange shut off before optimisation.

\chapter{TMCMC Implementation Details}
\label{app:tmcmc}

The two-phase GPU-accelerated TMCMC procedure of Chapter~\ref{ch:heine} is
summarised in Algorithm~\ref{alg:tmcmc}. The end-to-end GPU speedup ($\sim$$200\times$
on the matched benchmark) and production wall-clock are reported in
Table~\ref{tab:heine_timing}; the per-condition acceptance rates, log-evidence, and
fit metrics are in Table~\ref{tab:heine_rmse}.

\begin{algorithm}[htbp]
\caption{Two-phase GPU-accelerated TMCMC}
\label{alg:tmcmc}
\begin{algorithmic}[1]
\REQUIRE Data $\yobs$, viability data $\psi_i^{\mathrm{exp}}(t_k)$, $N_p$ particles, GPU device
\ENSURE MAP $\hat{\btheta}$, posterior samples, log-evidence $\ln\hat Z$
\STATE $\triangleright$ \textbf{Phase 1a: fix-$\psi$, wide prior}
\STATE Fix $\psi_i(t_k)\leftarrow\psi_i^{\mathrm{exp}}(t_k)$; draw $\{\btheta_j\}$ from the wide prior; $\beta_0\leftarrow0$
\WHILE{$\beta_m<1$}
  \STATE find $\beta_{m+1}$ s.t.\ $\mathrm{CoV}[\{w_j\}]=\delta_{\mathrm{target}}$ by bisection
  \STATE weights $w_j\leftarrow p(\yobs\mid\btheta_j)^{\beta_{m+1}-\beta_m}$; resample $\propto w_j$
  \STATE estimate weighted covariance $\hat{\boldsymbol\Sigma}^{(m)}$
  \FOR{$k=1$ \TO $K$}
    \STATE propose $\btheta_j'\sim\mathcal{N}(\btheta_j,\gamma^2\hat{\boldsymbol\Sigma}^{(m)})$ for all $j$
    \STATE evaluate all $N_p$ likelihoods in parallel via \texttt{jax.vmap} on GPU
    \STATE accept/reject (Metropolis--Hastings); clip to prior bounds
  \ENDFOR
\ENDWHILE
\STATE record posterior quantiles $q_{0.01},q_{0.99}$ and $\sigma$ per parameter
\STATE $\triangleright$ \textbf{Phase 1b: fix-$\psi$, narrowed prior} $[q_{0.01}-1.5\sigma,\,q_{0.99}+1.5\sigma]$; repeat
\STATE $\triangleright$ \textbf{Phase 2: free $\psi$, informative $\btheta$-prior from Phase 1b}
\STATE release $\psi_i$; augment likelihood with $-\sum_{k,i}(\psi_i(t_k)-\psi_i^{\mathrm{exp}}(t_k))^2/(2\sigma_\psi^2)$; repeat ($N_p=10{,}000$)
\RETURN $\hat{\btheta}$, $\hat{\bpsi}$, posterior samples, $\ln\hat Z$
\end{algorithmic}
\end{algorithm}

\paragraph{Hyperparameters.}
Adaptive tempering with CoV target $\delta_{\mathrm{target}}=1.0$; proposal scale
$\gamma^2=2.38^2/d$ adapted online toward acceptance $\bar\alpha\approx23\%$;
$K=20$ mutation steps per stage; forward model by implicit-Euler Newton
($\Delta t=10^{-4}$, $2{,}500$ steps), gauge $c^*=25$, barrier $K_p=10^{-4}$.
Convergence checked by $\mathrm{ESS}=N_p$ at the final stage and a two-seed
Gelman--Rubin criterion $\hat R<1.1$ on all active parameters.

\chapter{Spatial Reaction--Diffusion: Supplementary}
\label{app:pde}

\section*{PDE transport verification}

\begin{figure}[htbp]
  \centering
  \includegraphics[width=0.72\linewidth]{pde_verify_convergence.pdf}
  \caption[Diffusion-stencil convergence and boundary order]{Grid refinement of the diffusion
  stencil (reaction off, closed box). \textbf{Left:} error of the operator eigenvalue closest
  to the analytic $-D(\pi/L)^2$ Neumann mode; least-squares slope $2.00$ confirms
  $\mathcal{O}(\Delta z^2)$ interior accuracy. \textbf{Right:} node-$0$ truncation
  error at the ghost-node wall --- $\mathcal{O}(\Delta z)$ for a generic Neumann function.}
  \label{fig:pde_convergence}
\end{figure}

\begin{figure}[htbp]
  \centering
  \includegraphics[width=0.72\linewidth]{pde_verify_conservation.pdf}
  \caption[Mass conservation: diagonal vs.\ volume-filling cross-diffusion]{Closed-box
  conservation, diagonal (red) vs.\ volume-filling cross-diffusion (blue).
  \textbf{Top} (sub-critical, $\rho_{\max}\!\approx\!0.80$): diagonal drifts $+2.18\times10^{-3}$;
  cross-diffusion conserves to machine precision.
  \textbf{Bottom} (crowded, $\rho_{\max}\!=\!1.56$): diagonal clips exceed unity; cross-diffusion
  degeneracy caps $\rho$ near one.}
  \label{fig:pde_conservation}
\end{figure}

\begin{table}[htbp]
  \centering
  \caption[PDE transport verification results]{Measured verification outcomes
  (\texttt{verify\_pde\_numerics.py}, reaction off, closed box).}
  \label{tab:pde_verify}
  \small
  \begin{tabular}{llccc}
    \toprule
    Test & Scheme & Mass drift & $\rho_{\max}$ & Clip \\
    \midrule
    Convergence ($N_z\!=\!20$--$320$) & diagonal interior & \multicolumn{3}{l}{slope $2.00$ ($\mathcal{O}(\Delta z^2)$)}\\
    Boundary order (generic mode)     & ghost-node wall   & \multicolumn{3}{l}{slope $1.00$ ($\mathcal{O}(\Delta z)$)}\\
    \addlinespace
    Conservation, sub-critical & diagonal        & $+2.18\times10^{-3}$ & $0.79$ & $0$ \\
    Conservation, sub-critical & cross-diffusion & $0.0$ (mach.\ prec.) & $0.79$ & $0$ \\
    Conservation, crowded      & diagonal        & $+6.8\times10^{-4}$ & $1.56$ & $0$ \\
    Conservation, crowded      & cross-diffusion & $+1.9\times10^{-2}$ & $1.56$ & $2518$ \\
    \bottomrule
  \end{tabular}
\end{table}

\section*{FISH channel decode}

\begin{figure}[htbp]
  \centering
  \includegraphics[width=0.85\linewidth]{fig_fish_species_decode.png}
  \caption[HOBIC FISH: 5-species channel decode]{%
  Representative HOBIC FISH acquisition (commensal-HOBIC, Day~1, titanium substrate).
  Each column shows the confocal channel for one species
  (\textit{S.~oralis}, \textit{A.~naeslundii}, \textit{Vd/Vp},
  \textit{F.~nucleatum}, \textit{P.~gingivalis}) and the five-species composite
  (right); rows are independent fields of view.
  \textit{F.~nucleatum} is decoded from the blue $\cap$ red double-label channel.
  Scale bar: $50\,\mu\mathrm{m}$.}
  \label{fig:fish_species_decode}
\end{figure}

The species diffusivities $D_i$ are being calibrated against the CLSM $z$-profiles by a
grid/regularisation sweep on the cluster; the commensal-HOBIC fit has converged
(loss $\approx0.015$) while the dysbiotic-HOBIC fit is still in progress.

\chapter{In-Vivo Analysis: Supplementary Figures}
\label{app:dieckow_supp}

\begin{figure}[htbp]
  \centering
  \includegraphics[width=0.74\textwidth]{fig_loo_rmse_comparison.pdf}
  \caption[LOO-RMSE comparison across model variants]{LOO-RMSE ($\text{mean}\pm\text{std}$
  across 10 folds) for gLV and Hamilton variants. The Hamilton model with the full
  three-layer prior ($W{=}1.0$, $\bigstar$) matches the gLV free baseline ($0.0503$
  vs.\ $0.0501$) despite imposing $35$ metabolic sign constraints --- sign
  regularisation recovers predictive accuracy at no cost.}
  \label{fig:loo_rmse_comparison}
\end{figure}

\begin{figure}[htbp]
  \centering
  \includegraphics[width=0.70\textwidth]{dieckow_stability.pdf}
  \caption[LOO stability of the interaction matrix]{LOO stability (gLV
  $\alpha{=}0.25$, $10$ folds). \textbf{(A)} sign-consistency heatmap (boxed cells
  carry a sign constraint): $31/45$ at SC$=1.00$, $40/45$ at SC$\geq0.90$.
  \textbf{(B)} $|\bar A_{ij}|$ vs.\ fold-to-fold std by regime. \textbf{(C)}
  distributions of the strongest pairs.}
  \label{fig:dieckow_stability}
\end{figure}

\begin{figure}[htbp]
  \centering
  \includegraphics[width=0.62\textwidth]{dieckow_b_pca.pdf}
  \caption[Growth-vector clustering by community type]{PCA of patient growth
  vectors $\hat\bb^{(p)}$. PC1 separates CT1 (triangles) from CT2 (circles) with no
  label supervision --- patient-level ecology lives in $\bb$, not in the shared
  $\bA$.}
  \label{fig:dieckow_bpca}
\end{figure}
